\scanpage{10}\mspage{9}
en question sont dans $\unsure{\mathscr{L}|\Phi}$, i.e. que $\mathscr{L}$ est fermé dans $\mathscr{L}\cup\Psi$ (et itou pour le cas $A\in\Psi$, avec $\Psi$ fermé dans $\mathscr{L}\cup\Psi$). Donc il vient

\[
\begin{aligned}
2^{\circ})\quad &\deleted{\unsure{\Psi}}\mathscr{L}\ \text{fermé dans }\mathscr{L}\cup\Psi:\quad A\in\mathscr{L},\ B\in\Psi,\\
&B\subset\deleted{\unsure{B}}A\ \Longrightarrow\ B\in\deleted{B}\mathscr{L}.
\end{aligned}
\]
Mais on aura $\deleted{\unsure{\lvert B\rvert}}$ $B\subset\lvert\mathscr{L}\rvert\cap\lvert\Psi\rvert$, donc $B\in\mathscr{L}\Rightarrow B\deleted{\unsure{\subset}}\in\deleted{\unsure{\mathscr{L}}}\Psi$ OK.

\uline{\unsure{Définition|Definition}} On dit alors que $\mathscr{L}$ et $\Psi$ \uline{compatibles}

\uline{Corollaire \unsure{1|!}} Quand ces conditions sont satisfaites, $\lvert\deleted{A}\mathscr{L}\rvert\cap\lvert\Psi\rvert$ est une partie de $\lvert\mathscr{L}\rvert$ qui est $\mathscr{L}$-saturée, et une partie de $\Psi$ qui est $\Psi$-saturée, et les deux relations d'\unsure{équivalence|équilence} induites par $\lvert\deleted{\unsure{B}}\mathscr{L}\rvert$ et par $\lvert\Psi\rvert$ sont les $=\mathrm{s}$.

Mais attention, cette dernière condition ne suffit nullement à impliquer que $\mathscr{L}$ et $\Psi$ compatibles. Exemple : $\mathscr{L}$ et la famille \inserted{$\mathscr{L}^{\circ}$} des $A^{\circ}$ ont $=$ support, et définissent la $=$ relation d'équiv. sur $\lvert\mathscr{L}\rvert$, mais elles ne sont compatibles que si elles sont \uline{égales} \deleted{(car} i.e. $\mathscr{L}$ formé d'ens. deux à deux disjoints. En effet \deleted{(car la relation d'ordre sur $\mathscr{L}$ et sur $\mathscr{L}^{\circ}$ n'est la $=$ que si}

\uline{Corollaire 2} Supposons $\mathscr{L}$ et $\Psi$ compatibles. Alors $\lvert\mathscr{L}\rvert\subset\lvert\Psi\rvert$ sss $\mathscr{L}\subset\Psi$, $\lvert\mathscr{L}\rvert=\lvert\Psi\rvert$ sss $\mathscr{L}=\Psi$.
